{\parindent0pt
\begin{em}
% \blindtext
% \noindent
% El presente trabajo materializa el fin de un largo ciclo de formación profesional que comencé a recorrer allá por el año 2006, cuando debí abandonar mi ciudad natal, y junto a ella, a mi amada familia y amigos, para embarcarme en este largo viaje, impulsado por el viento de los sueños de un joven e ingenuo adolescente con hambre de aprender cada vez más y más.
El presente trabajo materializa el fin de un largo ciclo de formación académica que comencé a recorrer desde pequeño, cursando la tecnicatura en informática en el colegio N°20 de Villa Mercedes y que luego, en el año 2006, me llevó a abandonar mi ciudad natal, y junto a ella, a mi amada familia y amigos, para embarcarme en el largo viaje de la formación universitaria, impulsado por el fresco y vigoroso viento de los sueños de un joven e ingenuo adolescente ávido de conocimiento.

\

Agradezco a mi familia y antepasados; por ser parte de mi esencia ---trabajo, esfuerzo, sacrificio, valor, empuje, empatía, solidaridad, modestia, sencillez, austeridad.

\

Agradezco a mis amigos y seres queridos ---regocijo, júbilo, amparo, refugio, apoyo, reflexión, admiración, altruismo, generosidad.

\

Agradezco a los profesores del Colegio N°20, en especial a Luis Savid y a mi padre, por haber sabido despertar y estimular la llama de aquel joven.

\

Agradezco a la Universidad Nacional de San Luis por su excelencia, no sólo técnica sino también humana ---sus educadores e investigadores no sólo son sabios, sino también cálidos y humanos, con excelentes valores.

\

Por último, y no por ello menos importante, quiero agradecer a tres actores que tuvieron un rol clave en la materialización de este trabajo: a mi compañera de vida, Paula Yamanouchi, por el apoyo y el soporte cotidiano tanto en los momentos de estoicismo, como en los de aflicción y de euforia, a lo largo de estos años de doctorado; y a mis directores Marcelo Errecalde y Manuel Montes, quienes supieron guiarme tanto con la luz del conocimiento como la del corazón, con el consuelo, la paciencia, la empatía y los consejos propios de los de un padre a un hijo, estoy en deuda con ustedes.

\


\begin{flushright}
Eternamente agradecido por todo,\\
Sergio.
\end{flushright}

% \blindtext
\end{em}
}


